\section{Conclusion}\label{sec:conclusion}

Firms reallocate capacity as market opportunities change, and their decisions
alter the margins available to competitors.
Adjustment costs make these responses gradual, while investors can value the
anticipated cash-flow consequences when news arrives.
Building on \citet{xu_zenou_zhou_conflict_2022}, we characterize sequential
Nash capacity choices, establish equilibrium and comparative-static results,
and connect firms' operating paths to financial valuation.
In a restricted two-opportunity economy, firms with greater initial exposure
experience larger return responses to positive exogenous free cash flow shocks.
As another firm increases its exposure to the improving activity, its news
distribution moves closer to
the initial distributions of more-exposed competitors when reporting
intensity is equal across activities.
Under the stated conditions, this generates a nonnegative expected local
association between relative financial returns and changes in firms'
positions in information space.

We investigate this prediction using squared Wasserstein distances between
distributions of corporate-news embeddings and excess returns weighted by
article timing.
Monte Carlo experiments show that limited article coverage makes the positive
population relationship difficult to recover at the small
operating-adjustment scale considered.
Changes in the relative coverage of activities can also produce a positive
association while capacities remain fixed.
In a sample of \ppnum[0]{ir-firms} Nasdaq-listed firms over 2018--2022, the
estimated association is positive across four specifications but small.
A separate held-out exercise finds that the restricted movement pattern
accounts for little of the variation in annual news distributions.

By treating firms' economic footprints as outcomes of strategic choices, the
paper adds an explanation for how those footprints evolve to research on
distribution-valued firms.
The analysis also establishes why the production of financial news matters
for measuring those changes: operating reallocations and shifts in coverage
can have similar effects on observed distributions.
Independent, dated measures of operating activity would allow future research
to distinguish these mechanisms and assess the model more directly.
This provides a basis for studying changes in firms' economic footprints
through their reported activities and financial returns.
